\subsection{有限个测度的乘积}

设 \(X_{i}\) （ \(1 \leqslant i \leqslant n\) ）是 \(n\) 个局部紧空间， \(X = \prod_{i=1}^{n} X_{i}\) 是它们的乘积。形如

\[
\left(x _ {1}, x _ {2}, \dots , x _ {n}\right) \mapsto f _ {1} \left(x _ {1}\right) f _ {2} \left(x _ {2}\right) \dots f _ {n} \left(x _ {n}\right),
\]

（其中 \(f_{i} \in \mathcal{K}(\mathrm{X}_{i};\mathbf{C})\) ， \(1 \leqslant i \leqslant n\) ）的复函数的线性组合所成之集可典范地等同于张量积 \(\bigotimes_{i=1}^{n} \mathcal{K}(\mathrm{X}_{i};\mathbf{C})\) ，且由 No.~1 的引理 1（对 \(n\) 用归纳法）可知，这一张量积在 \(\mathcal{K}(\mathrm{X};\mathbf{C})\) 中稠密。

现在设 \(\mu_{i}\) 是 \(X_{i}\) （ \(1 \leqslant i \leqslant n\) ）上的一个测度；则在 X 上存在唯一的测度 \(\nu\) ，使得对 \(f_{i} \in \mathcal{K}(X_{i}; \mathbf{C})\) （ \(1 \leqslant i \leqslant n\) ），

\[
\left\langle f _ {1} \otimes f _ {2} \otimes \dots \otimes f _ {n}, \nu \right\rangle = \prod_ {i = 1} ^ {n} \left\langle f _ {i}, \mu_ {i} \right\rangle .\tag{11}
\]

因为若这一测度存在，则由前所述它是唯一的。另一方面，设 \(\nu = \mu_{1} \otimes \mu_{2} \otimes \cdots \otimes \mu_{n}\) 是由递归关系

\[
\mu_ {1} \otimes \mu_ {2} \otimes \dots \otimes \mu_ {n} = (\mu_ {1} \otimes \mu_ {2} \otimes \dots \otimes \mu_ {n - 1}) \otimes \mu_ {n}.
\]

在 X 上定义的测度。由 No.~1 的公式 (1) 及这一定义（对 n 用归纳法）可知 \(\nu\) 满足 (11)；称它为诸测度 \(\mu_{i}\) （ \(1 \leqslant i \leqslant n\) ）的乘积测度，并仍记作 \(\bigotimes_{i=1}^{n} \mu_{i}\) 。关系式 (11) 也可写成

\[
\left\langle f _ {1} \otimes f _ {2} \otimes \dots \otimes f _ {n}, \mu_ {1} \otimes \mu_ {2} \otimes \dots \otimes \mu_ {n} \right\rangle = \prod_ {i = 1} ^ {n} \left\langle f _ {i}, \mu_ {i} \right\rangle .\tag{12}
\]

命题 7（《测度乘积的结合律》）. --- 设 \((\mathbf{I}_k)_{1 \leqslant k \leqslant r}\) 是区间 \([1, n]\) （ \(\mathbf{N}\) 中的）的一个分割；则

\[
\bigotimes_ {k = 1} ^ {r} \left(\bigotimes_ {i \in I _ {k}} \mu_ {i}\right) = \bigotimes_ {i = 1} ^ {n} \mu_ {i}\tag{13}
\]

因为由 (12)，这两个测度对 \(\bigotimes_ {i = 1} ^ {n} \mathcal{K}(X_{i}; \mathbf{C})\) 中的每个函数取值相同。

函数 \(f \in \mathcal{K}(\mathbf{X};\mathbf{C})\) 关于乘积测度的积分记作

\[
\int f d \mu_ {1} d \mu_ {2} \dots d \mu_ {n},
\]

或

\[
\iint \ldots \int f d \mu_ {1} d \mu_ {2} \ldots d \mu_ {n}
\]

或

\[
\int f (\mu_ {1} \otimes \dots \otimes \mu_ {n})
\]

或

\[
\iint \dots \int f (x _ {1}, x _ {2}, \dots , x _ {n}) d \mu_ {1} (x _ {1}) d \mu_ {2} (x _ {2}) \dots d \mu_ {n} (x _ {n})
\]

或

\[
\iint \dots \int f (x _ {1}, x _ {2}, \dots , x _ {n}) \mu_ {1} (x _ {1}) \mu_ {2} (x _ {2}) \dots \mu_ {n} (x _ {n})
\]

（带 n 个积分号 \(\int\) ）；称它为 n 阶多重积分或 n 重积分。依据测度乘积的结合律及积分换序定理（No.~1, Th. 2），对每个置换 \(\sigma\) （ \([1,n]\) 的），有

\[
\iint \dots \int f d \mu_ {1} d \mu_ {2} \dots d \mu_ {n} = \int d \mu_ {\sigma (1)} \int d \mu_ {\sigma (2)} \dots \int f d \mu_ {\sigma (n)}.\tag{14}
\]

积分记号和公式 (14) 可以以明显的方式推广到函数 \(\mathbf{f} \in \mathcal{K}(\mathrm{X};\mathrm{E})\) （取值于 Hausdorff 局部凸空间 E 中、且 \(\mathbf{f}(\mathrm{X})\) 含于 E 的一个完备凸子集中的）。把 No. 2 与 No. 3 中关于两个测度的乘积的结果推广到任意有限个测度的乘积，这一工作留给读者。

特别地， \(\mathbf{R}^n\) 上的 Lebesgue 测度称为 \(n\) 个与 \(\mathbf{R}\) 上的 Lebesgue 测度相同的测度的乘积；满足前述条件的函数 \(\mathbf{f} \in \mathcal{K}(\mathbf{R}^n; \mathbf{E})\) 的积分记作

\[
\iint \dots \int \mathbf {f} (x _ {1}, x _ {2}, \dots , x _ {n}) d x _ {1} d x _ {2} \dots d x _ {n}
\]

并且等于

\[
\int_ {- \infty} ^ {+ \infty} d x _ {1} \int_ {- \infty} ^ {+ \infty} d x _ {2} \dots \int_ {- \infty} ^ {+ \infty} \mathbf {f} (x _ {1}, x _ {2}, \dots , x _ {n}) d x _ {n}.
\]

\(R^{n}\) 上的 Lebesgue 测度在每个平移下不变。
% semantic-fragments: legacy-input-seam
\relax{}
